\documentclass[11pt]{article}
\usepackage{amsmath, amssymb, amsthm}
\usepackage[margin=1.1in]{geometry}
\usepackage{parskip}
 
% Axler-style column notation: A_{.,k}
% We'll use a clean serif look consistent with math textbooks
 
\newtheorem*{claim}{Claim}
 
\title{\textbf{Verification: $AI = IA = A$ for Square Matrices}\\[0.4em]
       \large Following Axler, \textit{Linear Algebra Done Right 4th Ed.}, \S 3C--3D}
\date{}
 
\begin{document}
\maketitle
 
\begin{claim}
Let $A$ be an $n \times n$ matrix and let $I$ be the $n \times n$ identity matrix.
Then $AI = IA = A$.
\end{claim}
 
\medskip
 
We recall Theorem 3.51 from \S 3C: for two matrices $C$ (of size $m \times c$)
and $R$ (of size $c \times n$), the $k$th column of the product $CR$ is the linear
combination of the columns of $C$ with coefficients taken from the $k$th column
of $R$:
\[
    (CR)_{\cdot,\, k}
    \;=\;
    R_{1,k}\, C_{\cdot,\, 1}
    \;+\;
    R_{2,k}\, C_{\cdot,\, 2}
    \;+\;
    \cdots
    \;+\;
    R_{c,k}\, C_{\cdot,\, c}
    \;=\;
    \sum_{j=1}^{c} R_{j,k}\, C_{\cdot,\,j}.
\]
 
Also recall that the $k$th column of the identity matrix $I$ is the vector
$e_k$, which has a $1$ in row $k$ and $0$'s everywhere else:
\[
    I_{\cdot,\, k} \;=\; e_k, \qquad
    \text{so}\quad
    I_{j,\,k} =
    \begin{cases}
        1 & \text{if } j = k,\\
        0 & \text{if } j \neq k.
    \end{cases}
\]
 
\bigskip
 
\noindent\textbf{Part 1: $AI = A$.}
 
Fix any column index $k \in \{1, \dots, n\}$.
Applying Theorem 3.51(a) with $C = A$ and $R = I$:
\[
    (AI)_{\cdot,\, k}
    \;=\;
    \sum_{j=1}^{n} I_{j,\,k}\, A_{\cdot,\, j}.
\]
Since $I_{j,k} = 0$ for $j \neq k$ and $I_{k,k} = 1$, every term in the sum
vanishes except the $j = k$ term:
\[
    (AI)_{\cdot,\, k}
    \;=\;
    1 \cdot A_{\cdot,\, k}
    \;=\;
    A_{\cdot,\, k}.
\]
Because the $k$th column of $AI$ equals the $k$th column of $A$ for every $k$,
we conclude $AI = A$. \hfill$\square$
 
\bigskip
 
\noindent\textbf{Part 2: $IA = A$.}
 
Fix any column index $k \in \{1, \dots, n\}$.
Applying Theorem 3.51(a) with $C = I$ and $R = A$:
\[
    (IA)_{\cdot,\, k}
    \;=\;
    \sum_{j=1}^{n} A_{j,\,k}\, I_{\cdot,\, j}
    \;=\;
    A_{1,k}\, e_1
    \;+\;
    A_{2,k}\, e_2
    \;+\;
    \cdots
    \;+\;
    A_{n,k}\, e_n.
\]
Expanding each scaled standard basis vector, the $i$th entry of this sum is
\[
    \bigl[(IA)_{\cdot,\,k}\bigr]_{i}
    \;=\;
    \sum_{j=1}^{n} A_{j,k}\,(e_j)_i
    \;=\;
    A_{i,k},
\]
since $(e_j)_i = 1$ only when $j = i$ and is $0$ otherwise.
Hence $(IA)_{\cdot,\, k} = A_{\cdot,\, k}$ for every $k$, giving $IA = A$.
\hfill$\square$
 
\bigskip
 
\noindent\textbf{Conclusion.} Both $AI = A$ and $IA = A$, so $AI = IA = A$.
 
\end{document}